\noindent\textbf{ZUR NETZAUSGLEICHUNG.}\hfill319
\begin{center}\textbf{[10.]}\par\medskip
\textbf{[Briefwechsel zur Netzausgleichung.]}\end{center}

\begin{center}\textsc{Gauss} an \textsc{Olbers}.\quad Göttingen, 2. November 1823.\end{center}

.\ .\ .\ .\ .\ .\ Ich habe nunmehro die mühsame Ausgleichung meiner sämmtlichen
Messungen von 1821---1823, so weit sie die Hauptdreiecke betrifft,
vollendet, so dass nun nicht nur die Summen der Winkel der einzelnen Dreiecke,
sondern auch die Verhältnisse der Seiten in den gekreuzten Vierecken
und Fünfecken genau zu einander passen, und zwar nach den strengen Principien
der Wahrscheinlichkeitsrechnung sine ira et studio, und ohne alle
Willkürlichkeit. Es scheint nicht, dass man das letztere von den Messungen
anderer Geodäten sagen könne. Im ganzen Systeme sind 76 Richtungen, d.\ i.\
38 jede hinwärts und herwärts. Bei \textls[120]{keiner} von ihnen hat die Ausgleichung
$1^{\prime\prime}$ betragen; die grösste ist $0^{\prime\prime}{,}813$ bei der Richtung von Nindorf nach Hamburg,
wo das Pointiren auf den Michaelisthurm bei der heerrauchartigen Atmosphäre
und der Phasenstörung immer sehr schwierig war; die nächst grösste
ist $0^{\prime\prime}{,}788$ bei der Richtung von Lüneburg nach Wilsede, wo zwar der Zielpunkt
Heliotroplicht, aber die Aufstellung auf einem hölzernen Stativ in der
Laterne des Thurmes gewiss nicht von der Solidität war, wie auf Steinpostamenten,
und das Gewicht des Beobachters nach seiner verschiedenen
Stellung Einfluss auf das Instrument gehabt haben mag. Der \textls[120]{mittlere}
Fehler aller Richtungen, so verstanden wie in meiner Theoria Combinationis,
ist $0^{\prime\prime}{,}486$.

Es bilden sich in dem ganzen System 26 Dreiecke, in denen ich alle
Winkel gemessen habe. Der grösste Fehler der Summe der Winkel ist jetzt
$2^{\prime\prime}{,}175$ bei dem Dreieck Nindorf-Hamburg-Timpenberg, wo die Richtung von
Nindorf nach Hamburg vorzüglich Schuld haben mag. Der nächst grösste Fehler
ist, wie schon oben erwähnt, bei dem Dreiecke Brocken-Hohehagen-Hils; er
beträgt $1^{\prime\prime}{,}806$, und die Richtung vom Hils zum Brocken wird nun noch vorzüglich
Schuld sein.
\par\medskip
.\ .\ .\ .\ .\ .

\clearpage
\noindent320\hfill\textbf{BRIEFWECHSEL.}
\begin{center}\textsc{Gauss} an \textsc{Olbers}.\quad Zeven, 6. Julius 1824.\end{center}

\begingroup
\makeatletter
\renewcommand{\@makefnmark}{\hbox{\normalfont\textsuperscript{*})}}
\makeatother
.\ .\ .\ .\ .\ .\ Die neuen Messungen vom Jahr 1823 [haben] einige obwohl
sehr kleine Modificationen in dem ganzen System hervorgebracht\footnote{Auch ist in meiner jetzigen Rechnung \textsc{Walbecks} Abplattung \textls[120]{genau} zum Grunde gelegt, während bei der vorigen, wie Sie wissen, durch Versehen eine obwohl äusserst kleine Differenz statthatte.}.
\par\medskip
.\ .\ .\ .\ .\ .
\endgroup
